% ==============================================================================
% Fichier     : main.tex
% Exercice    : 3.4.3 — Première Mission
% Cours       : Analyse de Risque, Audit et Gestion de Crise
% Institution : ENSPY — Université de Yaoundé I
% ==============================================================================

\documentclass[12pt,a4paper]{article}
\usepackage{style_TCHADJIE}

% ==============================================================================
\begin{document}
% ==============================================================================

% -----------------------------------------------------------------------
% PAGE 1 — PAGE DE GARDE INSTITUTIONNELLE
% Tous les arguments sont passés ici, rien n'est hardcodé dans le corps
% -----------------------------------------------------------------------
\pagegardeENSP
    {ENSPY.png}
    {Analyse de Risque, Audit et Gestion de Crise}
    {EXERCICE 3.4.3}
    {22P092}
    {TCHADJIE TOUKAM LARRY}
    {}
    {MINKA MI NGUIDJOÏ Thierry Emmanuel}
    {2025/2026}

% -----------------------------------------------------------------------
% PAGE 2 — BANDEAU DE VERSION + NOTE LIMINAIRE + TABLE DES MATIÈRES
% -----------------------------------------------------------------------
\bandeauVersion{1.0}{Mars 2026}%
    {Soumis pour validation - COMEX LiZbiZ}

\vspace{0.6cm}

\begin{boxnote}
Ce document constitue notre proposition de charte de gouvernance simple pour l'entreprise LiZbiZ.
En tant que consultant externe fraîchement recruté par la DSI de LiZbiZ,
nous nous sommes appuyés sur les principes de bonne gouvernance du SI, les concepts RACI, les rôles RSSI/DSI/DPO et sur
les bonnes pratiques des chartes de gouvernance open source pour produire
ce livrable destiné au Comité de Direction.

Ce document est un \textbf{premier jet} : il sera mis à jour après les
retours du COMEX. Toute remarque du senior qui relira ce travail est
la bienvenue.
\end{boxnote}

\vspace{0.5cm}

\tableofcontents

\newpage

% ==============================================================================
\section{Contexte et problématique identifiée}
% ==============================================================================

\begin{boxmission}
LiZbiZ est une entreprise de e-commerce et logistique basée à Lyon,
spécialisée dans la vente de matériel informatique high-tech avec une
promesse de livraison en 24h. Après une croissance rapide (de 20 à
150 salariés en deux ans), l'entreprise fait face à un déficit structurel
de gouvernance de la sécurité. Cotée en bourse depuis six mois, elle
réalise un chiffre d'affaires de 45~M€. La mission nous a été confiée
par la DSI, Mme~Sophie~M., dans le contexte du recrutement d'un RSSI et
de la structuration de la gouvernance sécurité.
\end{boxmission}

\subsection{Lacunes actuelles documentées}

Avant de proposer une charte, il semble important de rappeler
brièvement l'état des lieux diagnostiqué :

\begin{boxcritique}
\begin{enumerate}
    \item \textbf{Conflit d'intérêt structurel.}
    La DSI (Mme~Sophie~M.) cumule les fonctions de DSI et de RSSI par
    défaut. Elle valide elle-même les droits d'accès qu'elle configure ---
    violation directe du principe de séparation des tâches.

    \item \textbf{Absence de comité de sécurité.}
    Aucune instance formelle ne remonte les risques à la Direction
    Générale. Le PDG, Jean-Pierre~L., n'est pas informé des
    vulnérabilités techniques (mises à jour irrégulières sur le serveur
    web, droits ERP trop larges pour des comptes d'ex-collaborateurs).

    \item \textbf{Pas de matrice RACI formalisée.}
    Les responsabilités de sécurité sont gérées «~à l'oral~», ce qui
    génère des oublis et des erreurs de processus, notamment lors de la
    gestion des accès Active Directory.

    \item \textbf{Poste de RSSI vacant.}
    La case \terme{Accountable} pour les décisions de sécurité est vide.
    C'est le risque organisationnel le plus immédiat : personne n'est
    habilité à décider seul en cas d'incident critique.
\end{enumerate}
\end{boxcritique}

\separateur

% ==============================================================================
\section{Charte de gouvernance de la sécurité --- LiZbiZ}
% ==============================================================================

\subsection{Objet et portée}

\begin{boxscope}
Cette charte s'applique à l'ensemble du Système d'Information de LiZbiZ,
y compris l'infrastructure cloud, les terminaux utilisateurs, les systèmes
logistiques IoT (entrepôt connecté), et les traitements de données à
caractère personnel soumis au \norme{RGPD}. Elle est opposable à tous les
employés, prestataires et partenaires ayant accès au SI de LiZbiZ.
\end{boxscope}

% -----------------------------------------------------------------------
\begin{boxarticle}{1 --- Principes fondamentaux de gouvernance}

\clause{Séparation des pouvoirs}
La \terme{Direction des Systèmes d'Information} (DSI) est chargée de
l'exploitation opérationnelle du SI. Le \terme{Responsable de la Sécurité
des Systèmes d'Information} (RSSI) est chargé du contrôle et de la
politique de sécurité. Ces deux fonctions sont \textbf{obligatoirement
exercées par des personnes distinctes} : la DSI ne peut pas être juge
et partie de sa propre sécurité.

\clause{Responsabilité ultime (Accountability)}
Conformément au principe défini dans le cours sous référence
\norme{COBIT}, la \terme{redevabilité} (Accountability) de la sécurité
du SI remonte jusqu'au Président Directeur Général. Le PDG répond devant
les actionnaires et les régulateurs en cas d'incident majeur. Cette
redevabilité est intransmissible.

\clause{Alignement stratégique}
Toute mesure de sécurité doit rester alignée avec les objectifs métier de
LiZbiZ. En particulier, les contrôles de sécurité ne doivent pas
compromettre la promesse de livraison en 24h qui constitue l'avantage
concurrentiel de l'entreprise \norme{ISO 27001}.

\clause{Cycle d'amélioration continue}
La gouvernance de la sécurité suit un cycle \terme{PDCA}
(Plan-Do-Check-Act) : la politique est planifiée, mise en œuvre, évaluée
annuellement, puis révisée. Tout incident de niveau C1
 déclenche une révision anticipée.

\end{boxarticle}

\vspace{0.4cm}

% -----------------------------------------------------------------------
\begin{boxarticle}{2 --- Rôle et missions du RSSI}

\begin{boxrole}{Fiche de poste synthétique : RSSI --- LiZbiZ}
\textbf{Rattachement hiérarchique :} Direction Générale (PDG) ---
\textit{pas} la DSI\\
\textbf{Statut actuel :} \statutVacant{} --- recrutement prioritaire\\
\textbf{Habilitation requise :} Accès aux données sensibles et aux
journaux de sécurité

\vspace{0.5em}
\textbf{\color{ChartePrimaire}Responsabilités principales :}
\begin{enumerate}
    \item \textbf{Définir et maintenir la PSSI.} Rédiger la Politique de
    Sécurité du SI en cohérence avec les contraintes réglementaires
    (\norme{RGPD}, \norme{PCI-DSS}) et les objectifs métier, puis en
    assurer la mise à jour annuelle.

    \item \textbf{Piloter l'analyse de risque.} Conduire l'analyse de
    risque selon la méthode \norme{EBIOS RM} et proposer des traitements
    au COMEX. Le RSSI évalue et communique sur le risque --- il n'est pas
    un technicien opérationnel.

    \item \textbf{Présider le Comité de Sécurité.} Animer la réunion
    mensuelle du Comité de Sécurité, qui rassemble DSI, DPO et
    représentants métiers.

    \item \textbf{Rendre compte au COMEX.} Présenter trimestriellement
    l'état de la sécurité pour validation des arbitrages budgétaires et
    des décisions de risque.

    \item \textbf{Gérer les incidents critiques.} Être l'autorité de
    décision lors des incidents de niveaux C1 et C2, et coordonner avec
    le DPO pour les obligations de notification réglementaire.

    \item \textbf{Sensibiliser les équipes.} Conduire des actions de
    formation à la sécurité pour l'ensemble des 150 collaborateurs.
\end{enumerate}

\vspace{0.4em}
\textbf{\color{ChartePrimaire}Décisions relevant du seul RSSI :}
\begin{itemize}
    \item Classification des incidents de sécurité (niveaux C1 à C4)
    \item Déclenchement de la procédure de réponse aux incidents
    \item Isolation ou cloisonnement d'urgence d'un système compromis
    \item Validation des exceptions à la politique de sécurité
\end{itemize}

\vspace{0.4em}
\textbf{\color{CharteCritBord}Ce que le RSSI ne fait \underline{pas} :}
\begin{itemize}
    \item Il ne configure pas directement les équipements (rôle DSI)
    \item Il ne valide pas seul un budget dépassant le seuil COMEX
    \item Il ne peut pas être subordonné à la DSI (cf. Article~1,
    clause~1)
\end{itemize}
\end{boxrole}

\clause{Rattachement et indépendance}
Le RSSI est rattaché \textbf{directement au PDG}, et non à la DSI.
Cette disposition garantit son indépendance de jugement. Une subordination
à la DSI recréerait exactement le conflit d'intérêt que cette charte
vise à corriger.

\clause{Période transitoire (avant recrutement)}
Dans l'attente du recrutement du RSSI, les attributions
\terme{Accountable} de la sécurité sont provisoirement assumées par le
PDG, assisté du consultant externe. La DSI conserve son rôle
d'exploitation mais \textbf{ne prend aucune décision de sécurité majeure
sans validation de la Direction Générale} pendant cette période.

\end{boxarticle}

\vspace{0.4cm}

% -----------------------------------------------------------------------
\begin{boxarticle}{3 --- Comité de Sécurité (CS) : instance opérationnelle}

\clause{Composition et fréquence}
Le Comité de Sécurité se réunit une fois par mois.

\vspace{0.4em}
\begin{center}
\begin{tabularx}{0.92\textwidth}{
    >{\bfseries\color{ChartePrimaire}}l
    >{\color{CharteArdoise}}X
    >{\centering\color{CharteSecondaire}\itshape\arraybackslash}p{2.4cm}}
\toprule
\enteteTableau{Membre} &
\enteteTableau{Rôle dans le Comité} &
\enteteTableau{Présence} \\
\midrule
\rowcolor{RACIPair}
RSSI (Président)  & Anime, présente le tableau de bord sécurité  & Obligatoire \\
\rowcolor{RACIImpair}
DSI               & Remonte les incidents techniques du mois      & Obligatoire \\
\rowcolor{RACIPair}
DPO               & Veille conformité RGPD et signalements CNIL   & Obligatoire \\
\rowcolor{RACIImpair}
Resp. Logistique  & Risques IoT et incidents entrepôt connecté    & Obligatoire \\
\rowcolor{RACIPair}
Représentant RH   & Incidents RH, arrivées et départs             & Sur invitation \\
\rowcolor{RACIImpair}
Auditeur interne  & Suivi des plans d'actions correctifs          & Trimestriel \\
\bottomrule
\end{tabularx}
\end{center}

\clause{Ordre du jour type}
Chaque séance aborde dans l'ordre : (1)~revue des incidents du mois,
(2)~suivi des plans d'actions correctifs, (3)~tableau de bord des risques,
(4)~points d'attention réglementaires, (5)~questions diverses.
Un compte rendu est rédigé et diffusé sous 48h.

\end{boxarticle}

\vspace{0.4cm}

% -----------------------------------------------------------------------
\begin{boxarticle}{4 --- Catégorisation des incidents critiques}
\label{sec:criticite}

\vspace{0.3em}
\begin{center}
\small
\begin{tabularx}{\textwidth}{
    >{\bfseries\centering\arraybackslash}p{1.2cm}
    >{\centering\arraybackslash}p{2.2cm}
    >{\color{CharteArdoise}}X
    >{\centering\color{CharteSecondaire}\arraybackslash}p{2.3cm}}
\toprule
\enteteTableau{Niveau} &
\enteteTableau{Qualification} &
\enteteTableau{Exemples LiZbiZ} &
\enteteTableau{Délai réponse} \\
\midrule
\rowcolor{CharteCritFond}
\textcolor{CharteCritBord}{C1} &
\textcolor{CharteCritBord}{Critique} &
Ransomware, fuite massive de données clients, compromission AD &
\textcolor{CharteCritBord}{\textbf{< 1h}} \\
\rowcolor{CharteRoleFond}
\textcolor{RACIAccountable}{C2} &
\textcolor{RACIAccountable}{Majeur} &
Intrusion confirmée, indisponibilité ERP $>$~2h, perte sauvegardes &
\textcolor{RACIAccountable}{\textbf{< 4h}} \\
\rowcolor{CharteActFond}
\textcolor{RACIConsulted}{C3} &
\textcolor{RACIConsulted}{Significatif} &
Tentative hameçonnage ciblée, anomalie de droits d'accès &
\textcolor{RACIConsulted}{\textbf{< 24h}} \\
\rowcolor{CharteNeutFond}
\textcolor{CharteGrisBleu}{C4} &
\textcolor{CharteGrisBleu}{Mineur} &
Alerte antivirus isolée, tentative d'accès échouée &
\textcolor{CharteGrisBleu}{\textbf{< 72h}} \\
\bottomrule
\end{tabularx}
\end{center}

\begin{boxnote}
Les critères combinent l'impact business (disruption livraison 24h) et
l'impact réglementaire (notification \norme{RGPD} sous 72h en cas de
violation de données personnelles).
\end{boxnote}

\end{boxarticle}

\newpage

% ==============================================================================
\section{Matrice RACI --- Gestion des incidents critiques}
% ==============================================================================

\subsection{Présentation de la matrice}

Cette matrice répond directement à la lacune identifiée qu'est : l'absence de RACI formel qui laisse les processus de sécurité se
gérer «~à l'oral~». Elle couvre les dix étapes du cycle de traitement
d'un incident, de la détection à la clôture formelle.

\begin{boxaction}
Cette matrice s'applique en priorité aux incidents \textbf{C1 et C2}.
La règle fondamentale a strictement respectée étant :
\textbf{une seule case A par ligne} (principe de l'unique tête de file
--- \textit{single point of accountability}).
\end{boxaction}

\vspace{0.4cm}

\subsection{La matrice}

\legendeRACI

\vspace{0.3em}
\begin{center}
\small\setlength{\tabcolsep}{3.5pt}
\begin{tabularx}{\textwidth}{
    >{\color{CharteArdoise}\raggedright\arraybackslash}p{4.9cm}
    >{\centering\arraybackslash}p{1.2cm}
    >{\centering\arraybackslash}p{1.2cm}
    >{\centering\arraybackslash}p{1.2cm}
    >{\centering\arraybackslash}p{1.2cm}
    >{\centering\arraybackslash}p{1.4cm}
    >{\centering\arraybackslash}p{1.4cm}}
\toprule
\enteteTableau{Activité du processus} &
\enteteTableau{RSSI} &
\enteteTableau{DSI} &
\enteteTableau{PDG} &
\enteteTableau{DPO} &
\enteteTableau{Resp. Logi.} &
\enteteTableau{Éq. Tech.} \\
\midrule
\rowcolor{RACIPair}
\textbf{1.} Détection et signalement
  & \raciI & \raciA & \raciI & \raciI & \raciI & \raciR \\
\rowcolor{RACIImpair}
\textbf{2.} Qualification et classification
  & \raciA & \raciR & \raciI & \raciC & \raciC & \raciR \\
\rowcolor{RACIPair}
\textbf{3.} Déclenchement procédure de réponse
  & \raciA & \raciR & \raciI & \raciI & \raciI & \raciI \\
\rowcolor{RACIImpair}
\textbf{4.} Isolation / confinement du système
  & \raciA & \raciC & \raciI & \raciVide & \raciC & \raciR \\
\rowcolor{RACIPair}
\textbf{5.} Communication de crise interne (COMEX)
  & \raciR & \raciC & \raciA & \raciI & \raciI & \raciVide \\
\rowcolor{RACIImpair}
\textbf{6.} Notification réglementaire RGPD / CNIL
  & \raciC & \raciI & \raciA & \raciR & \raciVide & \raciVide \\
\rowcolor{RACIPair}
\textbf{7.} Investigation forensique et collecte de preuves
  & \raciA & \raciC & \raciI & \raciI & \raciVide & \raciR \\
\rowcolor{RACIImpair}
\textbf{8.} Restauration des systèmes (PRA / PCA)
  & \raciC & \raciA & \raciI & \raciVide & \raciC & \raciR \\
\rowcolor{RACIPair}
\textbf{9.} Clôture formelle et documentation
  & \raciA & \raciR & \raciI & \raciC & \raciI & \raciR \\
\rowcolor{RACIImpair}
\textbf{10.} RETEX --- Analyse post-incident et plan correctif
  & \raciA & \raciR & \raciI & \raciC & \raciC & \raciR \\
\bottomrule
\end{tabularx}
\end{center}

\vspace{0.3em}
\legendeRACI

\subsection{Justification des choix clés}

\begin{enumerate}
    \item \textbf{Activité 1 --- Détection (DSI~=~A, Équipe~=~R).}
    L'équipe technique est en \textbf{R} car ce sont les ingénieurs qui
    voient les alertes en premier (SIEM, antivirus, journaux). La DSI
    est \textbf{A} car elle structure et maintient le dispositif de
    surveillance.

    \item \textbf{Activité 2 --- Qualification (RSSI~=~A).}
    Le RSSI maîtrise la grille de criticité. Confier ce rôle au PDG
    créerait un goulot d'étranglement et sortirait le PDG de son rôle
    stratégique.

    \item \textbf{Activité 5 --- Communication COMEX (PDG~=~A).}
    Le PDG décide de ce qui est communiqué vers les actionnaires et les
    médias. Le RSSI est \textbf{R} car il prépare les éléments techniques.

    \item \textbf{Activité 6 --- Notification RGPD (DPO~=~R, PDG~=~A).}
    Le DPO est \textbf{R} car la notification est sa mission
    réglementaire. Le PDG est \textbf{A} car elle engage juridiquement
    l'entreprise. Cette répartition respecte la distinction RSSI/DPO
    du cours.

    \item \textbf{Activité 8 --- Restauration (DSI~=~A).}
    La restauration technique relève du périmètre opérationnel de la DSI
    --- c'est le «~faire tourner le système~» au sens du cours (rôle
    \textit{Management} vs \textit{Gouvernance}).
\end{enumerate}

\separateur

% ==============================================================================
\section{Recommandations prioritaires}
% ==============================================================================

\begin{boxaction}
En tant que consultant, nous soumettons trois actions immédiates au
COMEX :

\begin{enumerate}
    \item \textbf{Recruter le RSSI (priorité absolue, horizon~: 3~mois).}
     Tant que ce poste reste vacant, la matrice RACI n’est pas utilisable.
    Le PDG a des tâches qui ne sont pas son cœur de métier,
    et le risque d’incident majeur sans responsable désigné est
    particulièrement élevé pour une entreprise cotée en bourse.

    \item \textbf{Réviser les droits d'accès ERP (quick win~: $<$~30
    jours).} Des comptes d'anciens collaborateurs restent actifs dans
    l'ERP Odoo --- violation du principe du moindre privilège
    \norme{ISO 27001} et vecteur de menace interne.

    \item \textbf{Formaliser un PCA minimal (horizon~: 6~mois).}
    Aucun Plan de Continuité d'Activité n'existe. En cas de ransomware,
    LiZbiZ n'a aucune procédure de bascule documentée --- risque
    existentiel pour un modèle fondé sur la disponibilité 24h/7j.
\end{enumerate}
\end{boxaction}

\newpage

% ==============================================================================
\section{Déclaration d'engagement et validation}
% ==============================================================================

\begin{boxarticle}{5 --- Entrée en vigueur et révision}

\clause{Entrée en vigueur}
Cette charte est applicable à compter de sa signature par le PDG et a
effet immédiat sur tous les acteurs du SI de LiZbiZ. 

\clause{Révision}
Cette charte est
revue annuellement lors de la revue de direction.
Tout incident de niveau C1 entraîne une révision anticipée
dans les 30~jours qui suivent la clôture.

\clause{Opposabilité}
Tout employé, prestataire ou partenaire qui a accès à notre SI est réputé avoir connaissance de la présente charte. Le non-respect de la présente charte est soumis à des sanctions disciplinaires.

\vspace{1.2cm}

\begin{center}
\renewcommand{\arraystretch}{2.2}
\begin{tabularx}{0.94\textwidth}{
    >{\centering\arraybackslash}X
    >{\centering\arraybackslash}X
    >{\centering\arraybackslash}X}
\hline
\rowcolor{RACIPair}
\textbf{\color{ChartePrimaire}PDG} &
\textbf{\color{ChartePrimaire}DSI} &
\textbf{\color{ChartePrimaire}RSSI} \\
Jean-Pierre L. & Sophie M. & \textit{Poste vacant} \\
\hline
\rowcolor{CharteNeutFond}
\footnotesize\textit{Signature \& Date} &
\footnotesize\textit{Signature \& Date} &
\footnotesize\textit{À pourvoir --- Recrutement en cours} \\
\hline
\end{tabularx}
\end{center}

\end{boxarticle}

\vspace{0.6cm}

% ==============================================================================
\end{document}
% ==============================================================================
